\conclusion
\begingroup\hyphenpenalty=10000\exhyphenpenalty=10000\sloppy
В рамках данной научно-исследовательской работы разработан и апробирован программный комплекс классификации Bitcoin-адресов по типу кошелька на основе датасета, сформированного в ходе НИРС-01 и НИРС-02. Реализован полный пайплайн подготовки данных: удаление коррелированных признаков, автоматический выбор стратегии устранения дисбаланса классов~(SMOTE превзошёл ADASYN по результатам прокси-сравнения на основе Random Forest) и стандартизация признаков. Обучено шесть моделей машинного обучения: MLP, Random Forest, XGBoost, LightGBM, Stacking и kNN. Для деревьевидных моделей произведена настройка гиперпараметров методом случайного поиска с двукратной стратифицированной кросс-валидацией. Проведена комплексная оценка качества: построены матрицы ошибок, ROC-кривые, вычислены взвешенная F1-мера и макро-усреднённый ROC-AUC для каждой модели. Сравнительный анализ показал явное преимущество деревьевидных методов над нейросетевой моделью на данной задаче: наилучший результат достигнут алгоритмом XGBoost~--- взвешенная F1-мера~0,7907, ROC-AUC~0,9610. Для лучшей модели выполнена оптимизация порогов классификации по критерию Юдена, повышающая полноту для миноритарных классов. Применение метода SHAP обеспечило интерпретируемость предсказаний и подтвердило значимость временны́х и экономических признаков. Детальный разбор ошибок выявил систематические паттерны смешения классов, прежде всего между \textit{services} и \textit{exchange}, и сформулированы гипотезы о путях улучшения. Практическая значимость работы определяется возможностью применения разработанного классификатора в системах AML-мониторинга криптовалютных операций, инструментах анализа блокчейн-активности и исследовательских задачах деанонимизации участников Bitcoin-сети. Результаты работы и исходный код доступны в открытом репозитории~\cite{github_repo_nirs3}. Перспективными направлениями дальнейших исследований являются: применение графовых нейронных сетей для учёта топологии транзакционного графа, расширение датасета и уточнение разметки класса~\textit{services}, разработка онлайн-системы обновления модели в условиях концептуального дрейфа, а также интеграция модели в реальную систему мониторинга.
\par\endgroup
